% 附录 A：拓扑与泛函分析预备知识（据原讲义第 57--66 页整理）
\chapter{拓扑与泛函分析预备知识}
\label{ch:appA}

\providecommand{\Xss}{\mathcal{X}^{**}}

本附录简要回顾正文用到的拓扑与泛函分析知识：点集拓扑的基本概念、拓扑向量空间、连续线性形式，
以及弱拓扑与弱*拓扑。最后证明（可分情形下的）Banach--Alaoglu 定理，并用它得到可分自反空间中单位球的弱紧性。
读者对前两节的内容大概已经熟悉，可以只浏览定义。

\section{点集拓扑}

\begin{definition}[拓扑]
\label{def:28}
集合 $X$ 上的一个\emph{拓扑}（topology）是 $X$ 的一族子集 $\tau$，其中的元素称为\emph{开集}（open set），满足：
\begin{enumerate}[label=(\roman*)]
  \item $\varnothing\in\tau$，$X\in\tau$；
  \item 有限个开集之交是开集：若 $\omega_1,\dots,\omega_k\in\tau$，则 $\omega_1\cap\dots\cap\omega_k\in\tau$；
  \item 任意多个开集之并是开集：若 $(\omega_i)_{i\in I}\subseteq\tau$，则 $\bigcup_{i\in I}\omega_i\in\tau$。
\end{enumerate}
此时称 $(X,\tau)$ 为\emph{拓扑空间}（topological space）。开集的补集称为\emph{闭集}（closed set）。
\end{definition}

\begin{example}[度量拓扑]
\label{exa:26}
设 $(\X,d)$ 是\emph{度量空间}（metric space），即 $d:\X\times\X\to\R$ 满足
\begin{enumerate}[label=(\roman*)]
  \item $d(x,y)\geqslant 0$，且 $d(x,y)=0$ 当且仅当 $x=y$；
  \item 对一切 $x,y\in\X$，$d(x,y)=d(y,x)$；
  \item 对一切 $x,y,z\in\X$，$d(x,z)\leqslant d(x,y)+d(y,z)$（三角不等式）。
\end{enumerate}
记闭球 $\mathrm{B}(x,r)=\{y\in\X\mid d(x,y)\leqslant r\}$。称 $\omega\subseteq\X$ 为开集，
如果对每个 $x\in\omega$，都存在 $r>0$ 使 $\mathrm{B}(x,r)\subseteq\omega$。
这样的开集全体 $\tau_d$ 满足拓扑的公理，称为由 $d$ 诱导的\emph{度量拓扑}（metric topology）。
\end{example}

\begin{definition}[邻域]
\label{def:29}
设 $(\X,\tau)$ 是拓扑空间，$x\in\X$。若集合 $N$ 包含某个含 $x$ 的开集，则称 $N$ 为 $x$ 的\emph{邻域}（neighborhood）。
$x$ 处的一个\emph{邻域基}（neighborhood basis）是 $\X$ 的一族子集 $\mathcal{B}_x$，满足：
(i) $\mathcal{B}_x$ 中每个集合都是 $x$ 的邻域；(ii) 对 $x$ 的任一邻域 $N$，存在 $B\in\mathcal{B}_x$ 使 $B\subseteq N$。
\end{definition}

\begin{example}
\label{exa:27}
在度量空间中，$\mathcal{B}_x=\{\mathrm{B}(x,\delta)\mid\delta>0\}$ 与 $\mathcal{B}'_x=\{\mathrm{B}(x,\tfrac1n)\mid n\in\N,\ n\geqslant 1\}$
都是 $x$ 处的邻域基。
\end{example}

\begin{definition}[一点处连续]
\label{def:30}
设 $f:(\X,\sigma)\to(\Y,\tau)$ 是拓扑空间之间的映射。若对 $f(x)$ 的任一邻域 $C$，原像 $f^{-1}(C)$ 都是 $x$ 的邻域，
则称 $f$ 在 $x$ 处\emph{连续}（continuous）。等价地，$f$ 在 $x$ 处连续当且仅当存在 $x$ 处的邻域基 $\mathcal{B}_x$
与 $f(x)$ 处的邻域基 $\mathcal{B}_{f(x)}$，使得
\[
  \forall C\in\mathcal{B}_{f(x)},\ \exists B\in\mathcal{B}_x,\quad f(B)\subseteq C\quad(\text{即 } B\subseteq f^{-1}(C)).
\]
对度量空间，取两边的邻域基为小球，就回到熟悉的 $\varepsilon$--$\delta$ 定义。
\end{definition}

连续性依赖于 $\sigma$ 与 $\tau$ 的选取。这一点要时刻记住，因为后面常在同一个空间上同时考虑两个拓扑。

\begin{definition}[连续]
\label{def:31}
若 $f:(\X,\sigma)\to(\Y,\tau)$ 在每个 $x\in\X$ 处连续，则称 $f$ 在 $\X$ 上连续。
等价地，开集的原像都是开集：$\forall\omega\in\tau$，$f^{-1}(\omega)\in\sigma$。
\end{definition}

\begin{definition}
\label{def:32}
称拓扑空间 $(\X,\tau)$ 是
\begin{itemize}
  \item \emph{分离的}或 \emph{Hausdorff} 的（separated / Hausdorff），如果任意两个不同点 $x\neq y$ 都有邻域 $O_x$、$O_y$ 使 $O_x\cap O_y=\varnothing$；
  \item \emph{可度量化的}（metrizable），如果 $\tau$ 是由某个度量诱导的。
\end{itemize}
\end{definition}

\begin{remark}
\label{rem:24}
度量空间自动是分离的。此外，在度量空间（从而在可度量化空间）中，一切拓扑概念都可以用序列刻画。例如：
\begin{itemize}
  \item 集合 $C\subseteq\X$ 是闭集，当且仅当 $C$ 中任一收敛序列的极限都属于 $C$；
  \item $f$ 在 $x$ 处连续，当且仅当对任一收敛到 $x$ 的序列 $(x_n)_{n\in\N}$，都有 $\lim_{n\to+\infty}f(x_n)=f(x)$。
\end{itemize}
对一般的拓扑空间（例如下面的弱拓扑），序列一般不足以刻画拓扑。
\end{remark}

\section{拓扑向量空间}

本课程的大部分结论在 $\X$ 为\emph{赋范空间}（normed space）时成立，即向量空间 $\X$ 配有范数 $\norm{\cdot}$，
其拓扑由度量 $d(x,y)=\norm{x-y}$ 诱导；此时以 $x$ 为心、$r$ 为半径的闭球记为
$\mathrm{B}(x,r)=\{y\in\X\mid\norm{x-y}\leqslant r\}$。特别地，这些结论适用于 Hilbert 空间、Banach 空间以及 $\R^d$。
不过，许多结论在更一般的（局部凸）拓扑向量空间中也成立；当需要考虑弱拓扑、弱*拓扑时，这种一般性正好派上用场。

\begin{definition}[拓扑向量空间]
\label{def:33}
\emph{拓扑向量空间}（topological vector space）是配有拓扑 $\tau$ 的向量空间 $\X$，使得向量加法
$(x,y)\in\X^2\mapsto x+y\in\X$ 与数乘 $(\lambda,x)\in\R\times\X\mapsto\lambda x\in\X$ 都连续。
\end{definition}

为了说清楚 $\X^2\to\X$、$\R\times\X\to\X$ 的连续性，需要在乘积空间上取定拓扑，即乘积拓扑。

\begin{definition}[乘积拓扑]
\label{def:34}
设 $(X_1,\tau_1)$、$(X_2,\tau_2)$ 是拓扑空间。称 $O\subseteq X_1\times X_2$ 在\emph{乘积拓扑}（product topology）$\tau_1\otimes\tau_2$ 下为开集，
如果对每个 $(x_1,x_2)\in O$，都存在 $x_1$ 的邻域 $N_1$ 与 $x_2$ 的邻域 $N_2$，使 $N_1\times N_2\subseteq O$。
\end{definition}

\paragraph{平移与数乘。}
在拓扑向量空间中，若 $O$ 是开集，则对任意 $x\in\X$，平移 $O+x=\{y+x\mid y\in O\}$ 也是开集。
这是因为映射 $T:y\mapsto y-x$ 连续（由加法的连续性），而 $O+x=T^{-1}(O)$。
同理，对 $\lambda\neq0$，$y\mapsto\lambda^{-1}y$ 连续，故 $\lambda O$ 是开集。
特别地，若 $\mathcal{B}$ 是原点处的邻域基，则 $\{B+x\mid B\in\mathcal{B}\}$ 是 $x$ 处的邻域基；
因此拓扑向量空间的拓扑完全由原点处的邻域决定。

\begin{definition}[局部凸拓扑向量空间]
\label{def:35}
若拓扑向量空间的原点有一个由凸集构成的邻域基，则称它是\emph{局部凸的}（locally convex）。
\end{definition}

\section{连续线性形式}

\begin{definition}[对偶空间]
\label{def:36}
$\X$ 上全体连续线性函数构成的空间称为 $\X$ 的\emph{拓扑对偶}（topological dual），记为 $\Xs$。
\end{definition}

$\Xs$ 的元素带星号记，如 $x^*\in\Xs$。为突出与 Hilbert 空间的类比，$\Xs$ 与 $\X$ 之间的\emph{对偶配对}（duality pairing）
用内积式的记号表示：$\pair{x^*}{x}\coloneqq x^*(x)$。

\begin{remark}
\label{rem:25}
当 $(\X,\norm{\cdot})$ 是赋范空间时，$\Xs$ 上可以赋予\emph{对偶范数}（dual norm）
\[
  \norm{x^*}_*=\sup_{\norm{x}\leqslant1}\pair{x^*}{x}.
\]
即使 $\X$ 不是赋范空间，也可以在 $\Xs$ 上构造一个“强”拓扑，但本课程用不到。
\end{remark}

下面的引理推广了赋范空间中“线性泛函连续当且仅当有界”的判据。

\begin{lemma}[连续线性形式的刻画]
\label{res:70}
设 $\phi:\X\to\R$ 是拓扑向量空间 $(\X,\sigma)$ 上的线性形式。以下几条等价：
\begin{enumerate}[label=(\roman*)]
  \item $\phi$ 连续；
  \item $\phi$ 在原点处连续，即对任意 $\varepsilon>0$，存在含原点的开集 $O$，在 $O$ 上 $\abs{\phi}<\varepsilon$；
  \item $\phi$ 在原点的某个邻域上有上界或有下界；
  \item $\phi$ 在某个非空开集上有上界或有下界。
\end{enumerate}
\end{lemma}

\begin{proof}
(i)$\Rightarrow$(ii) 与 (ii)$\Rightarrow$(iii) 是显然的（$\phi(0)=0$）。

(ii)$\Rightarrow$(i)：设 $x\in\X$，$\varepsilon>0$，取 (ii) 中的开集 $O$。则 $x+O$ 是 $x$ 的开邻域，且对 $z=x+h\in x+O$ 有
$\abs{\phi(z)-\phi(x)}=\abs{\phi(h)}<\varepsilon$。故 $\phi$ 在 $x$ 处连续。

(iii)$\Rightarrow$(iv)：邻域包含一个含原点的开集，在其上 $\phi$ 仍有同侧的界。

(iv)$\Rightarrow$(iii)：设 $U$ 是非空开集，$x_0\in U$，在 $U$ 上 $\phi\leqslant M$（有下界的情形同理）。
则 $U-x_0$ 是原点的开邻域，且对 $h\in U-x_0$，$\phi(h)=\phi(x_0+h)-\phi(x_0)\leqslant M-\phi(x_0)$。

(iii)$\Rightarrow$(ii)：不妨设在原点的某个邻域 $\omega$ 上 $\phi\leqslant M$（否则以 $-\phi$ 代替 $\phi$），并可设 $\omega$ 为开集、$M>0$。
由于 $x\mapsto-x$ 连续，$-\omega$ 也是原点的开邻域，于是 $\omega'\coloneqq\omega\cap(-\omega)$ 是原点的开邻域，并且对 $x\in\omega'$，
$\phi(x)\leqslant M$ 且 $-\phi(x)=\phi(-x)\leqslant M$，即 $\abs{\phi}\leqslant M$ 在 $\omega'$ 上成立。
对任意 $\varepsilon>0$，令 $\omega_\varepsilon=\frac{\varepsilon}{2M}\,\omega'$。它是含原点的开集，且在其上 $\abs{\phi}\leqslant\varepsilon/2<\varepsilon$。
这就证明了 $\phi$ 在原点处连续。
\end{proof}

\begin{lemma}[闭半空间]
\label{res:71}
设 $\phi:\X\to\R$ 是拓扑向量空间 $\X$ 上的线性形式，$\alpha\in\R$。以下几条等价：
\begin{enumerate}[label=(\roman*)]
  \item $\phi$ 连续；
  \item 半空间 $H=\{x\in\X\mid\phi(x)\leqslant\alpha\}$ 是闭集；
  \item 超平面 $H'=\{x\in\X\mid\phi(x)=\alpha\}$ 是闭集。
\end{enumerate}
\end{lemma}

\begin{proof}
若 $\phi$ 连续，则 $H$、$H'$ 分别是闭集 $(-\infty,\alpha]$、$\{\alpha\}$ 在 $\phi$ 下的原像，故为闭集。

(ii)$\Rightarrow$(i)：若 $H=\X$，则 $\phi$ 在 $\X$ 上有上界 $\alpha$；线性形式有上界只能是 $\phi=0$，自然连续。
否则 $\X\setminus H=\{\phi>\alpha\}$ 是非空开集，$\phi$ 在其上有下界 $\alpha$，由引理~\ref{res:70}(iv) 知 $\phi$ 连续。

(iii)$\Rightarrow$(i)：若 $\phi=0$，无需证明。设 $\phi\neq0$，则 $\phi$ 是满射，可取 $x_0$ 使 $\phi(x_0)<\alpha$。
注意 $\X\setminus H'=\{\phi\neq\alpha\}$ 是含 $x_0$ 的开集，但 $\phi$ 在其上未必有单侧界，需要再缩小。
$U\coloneqq(\X\setminus H')-x_0$ 是原点的开邻域。由数乘 $(\lambda,y)\mapsto\lambda y$ 在 $(0,0)$ 处的连续性，
存在 $\delta>0$ 与原点的开邻域 $W$，使得当 $\abs{\lambda}\leqslant\delta$、$w\in W$ 时 $\lambda w\in U$。
令 $V=\delta W$（开集，含原点）。则对 $y=\delta w\in V$ 及 $t\in[0,1]$，有 $ty=(t\delta)w\in U$，即 $x_0+ty\notin H'$。
于是 $t\mapsto\phi(x_0+ty)=\phi(x_0)+t\phi(y)$ 是 $[0,1]$ 上的连续函数，从小于 $\alpha$ 的值出发且始终不取值 $\alpha$，
由介值定理恒小于 $\alpha$；取 $t=1$ 得 $\phi(x_0+y)<\alpha$。
因此 $\phi$ 在非空开集 $x_0+V$ 上有上界 $\alpha$，由引理~\ref{res:70}(iv) 知 $\phi$ 连续。
\end{proof}

\section{弱拓扑与弱*拓扑}

很多情况下，范数拓扑“太细”：开集太多，紧集太少，难以证明最优化问题解的存在性。
因此我们希望定义更\emph{粗}（coarser）的拓扑，即开集更少、从而紧集更多的拓扑。
一个标准做法是：给定一族映射，构造使这些映射都连续的\emph{最粗}拓扑（按包含关系最小的拓扑）。
下面对向量空间 $\X$ 上的一族线性形式 $\Y$ 实施这一构造。

\subsection{由一族线性形式生成的拓扑}

\begin{lemma}
\label{res:72}
设 $\X$ 是向量空间，$\Y$ 是由 $\X$ 上的线性形式构成的向量空间。设 $\sigma$ 是 $\X$ 上的拓扑，且 $\Y$ 中每个线性形式关于 $\sigma$ 都连续。
则对任意 $x\in\X$、$y_1,\dots,y_N\in\Y$、$\varepsilon>0$，集合
\[
  B_{x,y_1,\dots,y_N,\varepsilon}\coloneqq\bigl\{z\in\X\bigm|\forall i\in\{1,\dots,N\},\ \abs{\pair{y_i}{z-x}}<\varepsilon\bigr\}
\]
关于 $\sigma$ 是开集。
\end{lemma}

\begin{proof}
由线性性，
\[
  B_{x,y_1,\dots,y_N,\varepsilon}=\bigcap_{i=1}^N y_i^{-1}\bigl((\pair{y_i}{x}-\varepsilon,\ \pair{y_i}{x}+\varepsilon)\bigr),
\]
它是开区间在 $\sigma$-连续映射 $y_1,\dots,y_N$ 下的原像的有限交，故为 $\sigma$-开集。
\end{proof}

这个引理引出下面的定义。

\begin{definition}[由线性形式生成的拓扑]
\label{def:37}
设 $\X$ 是向量空间，$\Y$ 是由 $\X$ 上的线性形式构成的向量空间。称 $O\subseteq\X$ 关于\emph{由 $\Y$ 生成的拓扑}
（topology generated by $\Y$）$\sigma(\X,\Y)$ 为开集，如果对每个 $x\in O$，存在 $y_1,\dots,y_N\in\Y$ 与 $\varepsilon>0$，
使 $B_{x,y_1,\dots,y_N,\varepsilon}\subseteq O$。
\end{definition}

集合 $B_{x,y_1,\dots,y_N,\varepsilon}$ 本身关于 $\sigma(\X,\Y)$ 是开集：若 $z\in B_{x,y_1,\dots,y_N,\varepsilon}$，令
$\delta=\varepsilon-\max_i\abs{\pair{y_i}{z-x}}>0$，则由三角不等式 $B_{z,y_1,\dots,y_N,\delta}\subseteq B_{x,y_1,\dots,y_N,\varepsilon}$。
因此 $x\in\X$ 处的一个邻域基为
\[
  \mathcal{B}_x=\{B_{x,y_1,\dots,y_N,\varepsilon}\mid N\in\N,\ y_1,\dots,y_N\in\Y,\ \varepsilon>0\}.
\]
在某种意义上，$B_{x,y_1,\dots,y_N,\varepsilon}$ 扮演了弱拓扑中“球”的角色。但这种拓扑有一个令人意外的特点：
$B_{x,y_1,\dots,y_N,\varepsilon}$ 包含仿射子空间 $x+\bigcap_{i=1}^N\ker y_i$，而当 $\X$ 无穷维时，有限个核之交（余维数至多为 $N$）
是非零子空间。所以若 $\X$ 是无穷维赋范空间，任何 $\sigma(\X,\Y)$-内部非空的集合都包含这样一个仿射子空间，因而无界。
特别地，无穷维时单位球 $\mathrm{B}(0,1)$ 的 $\sigma(\X,\Y)$-内部是空集。

下面的线性代数引理将在命题~\ref{res:73} 的证明中用到。

\begin{lemma}
\label{res:75}
设 $\X$ 是向量空间，$y_1,\dots,y_N,y$ 是 $\X$ 上的线性形式，满足 $\ker y_1\cap\dots\cap\ker y_N\subseteq\ker y$。
则 $y\in\spanop(y_1,\dots,y_N)$。
\end{lemma}

\begin{proof}
定义线性映射 $F:\X\to\R^{N+1}$，$F(x)=(y(x),y_1(x),\dots,y_N(x))$。
假设条件说明子空间 $L=F(\X)$ 不含点 $z=(1,0,\dots,0)$：若 $F(x)=z$，则 $x\in\bigcap_i\ker y_i$ 但 $y(x)=1$，矛盾。
记 $p\in L$ 为 $z$ 在 $L$ 上的正交投影，并记 $z-p=(\lambda,\lambda_1,\dots,\lambda_N)$，由 $z\notin L$ 知 $z-p\neq0$。
正交投影的刻画给出 $z-p\perp L$，即
\[
  \forall x\in\X,\quad \inner{F(x)}{z-p}=\lambda y(x)+\sum_{i=1}^N\lambda_iy_i(x)=0,
\]
也就是 $\lambda y+\sum_i\lambda_iy_i=0$。另一方面，$p\in L$ 蕴含 $\inner{p}{z-p}=0$，所以
\[
  \lambda=\inner{z}{z-p}=\inner{z-p}{z-p}+\inner{p}{z-p}=\norm{z-p}^2>0 .
\]
于是 $y=-\sum_i(\lambda_i/\lambda)\,y_i\in\spanop(y_1,\dots,y_N)$。
\end{proof}

\begin{proposition}
\label{res:73}
设 $\X$ 是向量空间，$\Y$ 是由 $\X$ 上的线性形式构成的向量空间。则
\begin{enumerate}[label=(\roman*)]
  \item $(\X,\sigma(\X,\Y))$ 是局部凸拓扑向量空间；
  \item 线性形式 $y:\X\to\R$ 关于 $\sigma(\X,\Y)$ 连续，当且仅当 $y\in\Y$。
\end{enumerate}
\end{proposition}

\begin{proof}
(i) \emph{$\sigma(\X,\Y)$ 是拓扑}。显然 $\varnothing,\X\in\sigma(\X,\Y)$，且定义中的条件对任意并封闭。
对有限交：若 $x\in O\cap O'$，$B_{x,y_1,\dots,y_N,\varepsilon}\subseteq O$，$B_{x,y'_1,\dots,y'_M,\varepsilon'}\subseteq O'$，则
\[
  B_{x,y_1,\dots,y_N,y'_1,\dots,y'_M,\min(\varepsilon,\varepsilon')}\subseteq B_{x,y_1,\dots,y_N,\varepsilon}\cap B_{x,y'_1,\dots,y'_M,\varepsilon'}\subseteq O\cap O'.
\]

\emph{加法连续}。记 $S(x,y)=x+y$。取 $x,y\in\X$，$z=x+y$，以及 $z$ 的任一邻域 $N_z$；它包含某个
$B_z\coloneqq B_{z,y_1,\dots,y_N,\varepsilon}$。令 $B_x=B_{x,y_1,\dots,y_N,\varepsilon/2}$，$B_y=B_{y,y_1,\dots,y_N,\varepsilon/2}$。
对 $x'\in B_x$、$y'\in B_y$，$z'=x'+y'$ 满足
\[
  \abs{\pair{y_i}{z'-z}}=\abs{\pair{y_i}{x'-x}+\pair{y_i}{y'-y}}\leqslant\abs{\pair{y_i}{x'-x}}+\abs{\pair{y_i}{y'-y}}<\varepsilon,
\]
故 $S(B_x\times B_y)\subseteq B_z\subseteq N_z$，而 $B_x\times B_y$ 是 $(x,y)$ 在乘积拓扑下的邻域。所以 $S$ 在 $(x,y)$ 处连续。

\emph{数乘连续}。记 $P(\lambda,x)=\lambda x$，取 $(\lambda,x)$ 及 $\lambda x$ 的邻域，它包含某个 $B_{\lambda x,y_1,\dots,y_N,\varepsilon}$。
令 $c=\max_i\abs{\pair{y_i}{x}}$，取 $\delta\in(0,1]$ 使 $(\abs{\lambda}+1+c)\delta<\varepsilon$。对 $\abs{\lambda'-\lambda}<\delta$、$x'\in B_{x,y_1,\dots,y_N,\delta}$，
\[
  \abs{\pair{y_i}{\lambda'x'-\lambda x}}
  =\abs{\lambda'\pair{y_i}{x'-x}+(\lambda'-\lambda)\pair{y_i}{x}}
  \leqslant(\abs{\lambda}+1)\delta+\delta c<\varepsilon,
\]
所以 $P$ 在 $(\lambda,x)$ 处连续。于是 $(\X,\sigma(\X,\Y))$ 是拓扑向量空间。
最后，原点的邻域基 $\{B_{0,y_1,\dots,y_N,\varepsilon}\}$ 由凸集组成（它们是有限个“带形”$\{\abs{y_i}<\varepsilon\}$ 之交），故空间局部凸。

(ii) 若 $y\in\Y$，则对任意开区间 $(a,b)$ 与 $x\in y^{-1}((a,b))$，取 $\varepsilon=\min(y(x)-a,\,b-y(x))$，
有 $B_{x,y,\varepsilon}\subseteq y^{-1}((a,b))$；故 $y^{-1}((a,b))$ 是开集，$y$ 连续。

反之，设线性形式 $y$ 关于 $\sigma(\X,\Y)$ 连续。由于 $y(0)=0$，存在原点的邻域 $N$，在其上 $\abs{y}<1$。
按拓扑的构造，存在 $y_1,\dots,y_N\in\Y$ 与 $\varepsilon>0$，使 $B_{0,y_1,\dots,y_N,\varepsilon}\subseteq N$，即
\[
  \bigl(\forall i\in\{1,\dots,N\},\ \abs{\pair{y_i}{x}}<\varepsilon\bigr)\ \Longrightarrow\ \abs{\pair{y}{x}}<1 .
\]
若 $x\in\ker y_1\cap\dots\cap\ker y_N$，则对任意 $t\in\R$，$tx$ 满足左边的条件，所以 $\abs{t\,y(x)}<1$ 对一切 $t$ 成立，只能 $y(x)=0$。
这说明 $\ker y_1\cap\dots\cap\ker y_N\subseteq\ker y$，由引理~\ref{res:75}，$y\in\spanop(y_1,\dots,y_N)\subseteq\Y$。
\end{proof}

\begin{corollary}
\label{res:74}
设 $y:\X\to\R$ 是线性形式，$\alpha\in\R$。则超平面 $y^{-1}(\alpha)$（相应地，半空间 $y^{-1}((-\infty,\alpha])$）关于 $\sigma(\X,\Y)$ 是闭集，
当且仅当 $y\in\Y$。
\end{corollary}

\begin{proof}
由命题~\ref{res:73}(i)，$(\X,\sigma(\X,\Y))$ 是拓扑向量空间，故可对它应用引理~\ref{res:71}：
超平面（半空间）闭当且仅当 $y$ 关于 $\sigma(\X,\Y)$ 连续，再由命题~\ref{res:73}(ii) 即当且仅当 $y\in\Y$。
\end{proof}

\begin{remark}[最粗拓扑]
\label{rem:26}
引理~\ref{res:72} 表明：若 $\sigma$ 是 $\X$ 上使 $\Y$ 中线性形式全都连续的拓扑，则 $\sigma$ 包含所有形如 $B_{x,y_1,\dots,y_N,\varepsilon}$ 的集合，
从而包含它们的任意并，即 $\sigma\supseteq\sigma(\X,\Y)$。结合命题~\ref{res:73}(ii)，$\sigma(\X,\Y)$ 是 $\X$ 上使 $\Y$ 中线性形式全都连续的\emph{最粗}拓扑。
（这也意味着它拥有最多的紧集、最多的收敛序列。）
\end{remark}

\begin{remark}[序列的收敛]
\label{rem:27}
$\X$ 中的序列 $(x_n)_{n\geqslant1}$ 关于 $\sigma(\X,\Y)$ 收敛到 $x\in\X$（即对 $x$ 的任一邻域 $O$，当 $n$ 充分大时 $x_n\in O$），当且仅当
\[
  \forall y\in\Y,\quad\lim_{n\to+\infty}\pair{y}{x_n}=\pair{y}{x}.
\]
事实上，必要性来自每个 $y\in\Y$ 的连续性；充分性：$x$ 的任一邻域 $O$ 包含某个 $B_{x,y_1,\dots,y_N,\varepsilon}$，
而由假设当 $n$ 充分大时对 $i=1,\dots,N$ 同时有 $\abs{\pair{y_i}{x_n-x}}<\varepsilon$，即 $x_n\in O$。
\end{remark}

\subsection{弱拓扑}

\begin{definition}[弱拓扑]
\label{def:38}
设 $\X$ 是拓扑向量空间，$\Xs$ 是 $\X$ 上连续线性形式构成的空间。由 $\Xs$ 生成的拓扑 $\sigma(\X,\Xs)$ 称为 $\X$ 上的\emph{弱拓扑}（weak topology）。
\end{definition}

\begin{remark}[与强拓扑的关系]
\label{rem:28}
为与弱拓扑 $\sigma(\X,\Xs)$ 区分，称 $\X$ 原有的拓扑为\emph{强拓扑}（strong topology）。
构成弱拓扑基的集合 $B_{x,x^*_1,\dots,x^*_N,\varepsilon}$ 关于原拓扑是开集（引理~\ref{res:72}），因此弱开集都是强开集，弱闭集都是强闭集。
也许有些反直觉的是：若函数 $f:\X\to\R$ 弱连续，则它强连续。（证明：对 $\R$ 的任一开集 $O$，$f^{-1}(O)$ 弱开，从而强开。）
反方向的蕴含一般不成立，但加上凸性假设后成立（例如见命题~\ref{res:12}）。
\end{remark}

\subsection{弱*拓扑}

\begin{definition}[弱*拓扑]
\label{def:39}
设 $\X$ 是拓扑向量空间，考虑\emph{典范映射}（canonical injection）
\[
  i:\X\to\Xss,\qquad x\mapsto\bigl(x^*\mapsto\pair{x^*}{x}\bigr).
\]
借助它，可以把 $\X$ 看作 $\Xs$ 上的一族线性形式 $i(\X)$。$\Xs$ 上由 $i(\X)$ 生成的拓扑称为\emph{弱*拓扑}（weak* topology），记为
$\sigma(\Xs,\X)\coloneqq\sigma(\Xs,i(\X))$。
\end{definition}

关于弱拓扑与弱*拓扑更系统的讨论可参见 [Bre10]（Brezis 的泛函分析教材）。

\begin{remark}[弱*收敛与逐点收敛]
\label{rem:29}
由注~\ref{rem:27}（取 $\Y=i(\X)$），$\Xs$ 中的序列 $(x^*_n)$ 关于 $\sigma(\Xs,\X)$ 收敛到 $x^*$（简称 $(x^*_n)$ \emph{弱*收敛}（weak-* converges）到 $x^*$）当且仅当
\[
  \forall x\in\X,\quad\lim_{n\to+\infty}\pair{x^*_n}{x}=\pair{x^*}{x}.
\]
换言之，把这些线性形式看作 $\X$ 上的函数，弱*收敛就是\emph{逐点收敛}（pointwise convergence）。
\end{remark}

\begin{remark}[分离性]
\label{rem:30}
$(\Xs,\sigma(\Xs,\X))$ 是分离的（Hausdorff）。事实上，设 $x^*\neq y^*$，则存在 $x\in\X$ 使 $\pair{x^*}{x}\neq\pair{y^*}{x}$，不妨设 $\pair{x^*}{x}<\pair{y^*}{x}$。
令 $r=\pair{\tfrac12(x^*+y^*)}{x}$，则两个弱*开集 $O_-=\{\xi\in\Xs\mid\pair{\xi}{x}<r\}$ 与 $O_+=\{\xi\in\Xs\mid\pair{\xi}{x}>r\}$ 互不相交，且 $x^*\in O_-$、$y^*\in O_+$。
\end{remark}

\begin{example}[测度]
\label{exa:28}
设 $K\subseteq\R^d$ 紧，$\X=\mathcal{C}^0(K)$ 是 $K$ 上连续函数的空间，配以一致范数 $\norm{\cdot}_\infty$。$\X$ 上有哪些连续线性泛函？
若 $\mu$ 是 $K$ 上的有限（带符号）Radon 测度，则 $f\in\mathcal{C}^0(K)\mapsto\int f\,\mathrm{d}\mu$ 是连续线性泛函。
反之，\emph{Riesz--Markov 定理}断言 $\X$ 上的每个连续线性泛函都由某个有限带符号 Radon 测度给出。因此本课程把 Radon 测度空间\emph{定义}为
\[
  \mathcal{M}(K)\coloneqq\mathcal{C}^0(K)^* .
\]
Radon 测度的一个例子是 \emph{Dirac 测度}（Dirac mass）$\delta_x\in\mathcal{M}(K)$：
\[
  \forall\phi\in\mathcal{C}^0(K),\quad\pair{\delta_x}{\phi}=\delta_x(\phi)\coloneqq\phi(x).
\]
$\mathcal{M}(K)$ 上的对偶范数称为测度的\emph{全变差}（total variation）：$\norm{\nu}_{TV}=\sup\{\pair{\nu}{\phi}\mid\norm{\phi}_\infty\leqslant1\}$。
若 $x\neq y$ 是 $K$ 中两点，可以构造 $\phi\in\mathcal{C}^0(K)$ 使 $\phi(x)=1$、$\phi(y)=-1$、$\norm{\phi}_\infty\leqslant1$，因此
\[
  \norm{\delta_x-\delta_y}_{TV}\geqslant\int\phi\,\mathrm{d}(\delta_x-\delta_y)=\phi(x)-\phi(y)=2 .
\]
于是，若 $x_n\to x$ 且 $x_n\neq x$，Dirac 测度 $\delta_{x_n}$ 关于全变差范数\emph{不}收敛到 $\delta_x$。另一方面，
\[
  \forall\phi\in\mathcal{C}^0(K),\quad\pair{\delta_{x_n}}{\phi}=\phi(x_n)\xrightarrow[n\to+\infty]{}\phi(x)=\pair{\delta_x}{\phi},
\]
所以 $(\delta_{x_n})$ 弱*收敛到 $\delta_x$，即关于 $\sigma(\mathcal{M}(K),\mathcal{C}^0(K))$ 收敛。
这说明弱*拓扑 $\sigma(\mathcal{M}(K),\mathcal{C}^0(K))$ 拥有更多的收敛序列，因而也有更多的紧集。
\end{example}

\begin{remark}[$(\Xs,\sigma(\Xs,\X))$ 上的连续线性形式]
\label{rem:31}
由命题~\ref{res:73}(ii)（取 $\Y=i(\X)$），$\Xs$ 上的线性形式关于弱*拓扑连续，当且仅当它由 $\X$ 的某个元素给出；
即 $(\Xs,\sigma(\Xs,\X))$ 的对偶可以等同于 $\X$。例如，配以弱*拓扑的测度空间的对偶是
$(\mathcal{M}(K),\sigma(\mathcal{M}(K),\mathcal{C}^0(K)))^*\simeq\mathcal{C}^0(K)$，而 $(\mathcal{M}(K),\norm{\cdot}_{TV})^*$ 则是一个大得多、
也难以描述得多的空间（有专著讨论这一问题，[Kap11]）。
\end{remark}

\subsection{Banach--Alaoglu 定理}

\begin{definition}[可分性]
\label{def:40}
若拓扑空间 $\X$ 含有可数稠密子集，则称它是\emph{可分的}（separable）。
\end{definition}

\begin{example}
\label{exa:29}
可分空间的例子：
\begin{itemize}
  \item 紧集 $K\subseteq\R^d$ 上的连续函数空间 $(\mathcal{C}^0(K),\norm{\cdot}_\infty)$。这是 Stone--Weierstrass 定理的推论：
  紧集上的连续函数可被多项式一致逼近，而有理系数多项式只有可数个；
  \item $p\in[1,+\infty)$ 时，$p$ 次可和序列空间 $(\ell^p,\norm{\cdot}_p)$，以及 $\R^d$ 中有界区域 $\Omega$ 上 $p$ 次可积函数空间 $(\mathrm{L}^p(\Omega),\norm{\cdot}_p)$。
  注意 $\ell^\infty$ 与 $\mathrm{L}^\infty(\Omega)$ \emph{不}可分。
\end{itemize}
\end{example}

\begin{theorem}[对偶空间中有界集的可度量化]
\label{res:76}
设 $\X\neq\{0\}$ 是可分赋范空间，$(x_n)_{n\geqslant1}$ 是 $\X$ 中由非零向量组成的稠密序列。
则在 $\Xs$ 的任一有界子集 $S$ 上，弱*拓扑（的限制）由下面的距离诱导：
\[
  d(x^*,y^*)=\sum_{n\geqslant1}2^{-n}\Bigl|\Bigl\langle x^*-y^*\Bigm|\frac{x_n}{\norm{x_n}}\Bigr\rangle\Bigr| .
\]
\end{theorem}

\begin{remark}
\label{rem:32}
这里 $S$ 上的弱*拓扑指子空间拓扑：$S$ 的弱*开集就是形如 $O'=O\cap S$ 的集合，其中 $O$ 是 $\Xs$ 的弱*开集。
\end{remark}

\begin{proof}
记 $e_n=x_n/\norm{x_n}$，$\norm{e_n}=1$。

\emph{$d$ 是 $\Xs$ 上的距离}。由对偶范数的定义，$\abs{\pair{x^*-y^*}{e_n}}\leqslant\norm{x^*-y^*}_*$，故级数收敛且 $d(x^*,y^*)\leqslant\norm{x^*-y^*}_*$。
对称性与三角不等式是显然的。若 $d(x^*,y^*)=0$，则对一切 $n$ 有 $\pair{x^*}{x_n}=\pair{y^*}{x_n}$；
两个连续线性形式在稠密子集上相等，故在整个 $\X$ 上相等，即 $x^*=y^*$。

设 $S\subseteq\mathrm{B}(0,R)$，$R>0$。于是对 $y^*,z^*\in S$，$\norm{y^*-z^*}_*\leqslant2R$。

\emph{$d$-开集是弱*开集}。设 $O'$ 是 $S$ 中关于 $d$ 的开集，$x^*\in O'$。存在 $\varepsilon>0$ 使 $\{z^*\in S\mid d(x^*,z^*)<\varepsilon\}\subseteq O'$。
取 $N\geqslant2$ 使 $2R\sum_{n\geqslant N}2^{-n}=2R\cdot2^{1-N}\leqslant\varepsilon/2$，则
\[
  \forall z^*\in S,\quad\sum_{n\geqslant N}2^{-n}\abs{\pair{x^*-z^*}{e_n}}\leqslant\frac{\varepsilon}{2}.
\]
若 $z^*\in S\cap B_{x^*,e_1,\dots,e_{N-1},\varepsilon/4}$（这是 $x^*$ 的弱*基本邻域与 $S$ 之交），则对 $1\leqslant n<N$ 有 $\abs{\pair{x^*-z^*}{e_n}}<\varepsilon/4$，从而
\[
  d(x^*,z^*)=\sum_{1\leqslant n<N}2^{-n}\abs{\pair{x^*-z^*}{e_n}}+\sum_{n\geqslant N}2^{-n}\abs{\pair{x^*-z^*}{e_n}}<\frac{\varepsilon}{4}+\frac{\varepsilon}{2}<\varepsilon .
\]
于是 $S\cap B_{x^*,e_1,\dots,e_{N-1},\varepsilon/4}\subseteq O'$。对每个 $x^*\in O'$ 取这样的弱*开集并求并，得 $O'=O\cap S$，其中 $O$ 是弱*开集，即 $O'$ 在 $S$ 中弱*开。

\emph{弱*开集是 $d$-开集}。设 $O'=O\cap S$，$O$ 弱*开，$x^*\in O'$。存在 $z_1,\dots,z_k\in\X$ 与 $\varepsilon>0$ 使 $B_{x^*,z_1,\dots,z_k,\varepsilon}\subseteq O$。
由稠密性，对每个 $j$ 取 $n_j$ 使 $\norm{z_j-x_{n_j}}<\varepsilon/(4R)$。对 $y^*\in S$，
\[
  \abs{\pair{x^*-y^*}{z_j}}\leqslant\abs{\pair{x^*-y^*}{x_{n_j}}}+\norm{x^*-y^*}_*\norm{z_j-x_{n_j}}
  <2^{n_j}\norm{x_{n_j}}\,d(x^*,y^*)+\frac{\varepsilon}{2},
\]
这里用到了 $2^{-n_j}\abs{\pair{x^*-y^*}{e_{n_j}}}\leqslant d(x^*,y^*)$。取 $\delta=\min_j\varepsilon/(2^{n_j+1}\norm{x_{n_j}})$，
则 $y^*\in S$、$d(x^*,y^*)<\delta$ 蕴含 $y^*\in B_{x^*,z_1,\dots,z_k,\varepsilon}\cap S\subseteq O'$。故 $O'$ 是 $d$-开集。
\end{proof}

\begin{theorem}[Banach--Alaoglu，序列形式]
\label{res:77}
设 $\X$ 是可分赋范空间。则 $\Xs$ 中任一有界序列都有弱*收敛的子序列。
\end{theorem}

\begin{proof}
设 $A=\{a_1,a_2,\dots\}$ 是 $\X$ 的可数稠密子集，$(x^*_n)$ 是 $\Xs$ 中的有界序列，$\norm{x^*_n}_*\leqslant R$（$R>0$）。

\emph{对角线论证}。对每个 $a\in A$，实数列 $(\pair{x^*_n}{a})_n$ 有界（$\leqslant R\norm{a}$），因而有收敛子列。
先在 $a_1$ 处抽取收敛子列，再从中在 $a_2$ 处抽取收敛子列，依此类推；取第 $k$ 个子列的第 $k$ 项组成对角子列。
它在每个 $a_k$ 处都收敛（从第 $k$ 项起它是第 $k$ 个子列的子列）。重新编号后，不妨设对每个 $a\in A$，存在 $f_a\in\R$ 使
$\lim_{n\to+\infty}\pair{x^*_n}{a}=f_a$。

\emph{在整个 $\X$ 上逐点收敛}。对一切 $n$ 及 $x,y\in\X$，$\abs{\pair{x^*_n}{x}-\pair{x^*_n}{y}}\leqslant R\norm{x-y}$。
令 $n\to+\infty$ 知 $f:a\mapsto f_a$ 在稠密集 $A$ 上是 $R$-Lipschitz 的，故可唯一延拓为 $\X$ 上的 $R$-Lipschitz 函数 $\hat f$。
下证 $x^*_n$ 在 $\X$ 上逐点收敛到 $\hat f$。设 $x\in\X$，$\varepsilon>0$。由稠密性取 $y\in A$ 使 $\norm{x-y}\leqslant\frac{\varepsilon}{3R}$；
由 $x^*_n(y)\to\hat f(y)$，存在 $N$ 使 $n\geqslant N$ 时 $\abs{\hat f(y)-x^*_n(y)}\leqslant\frac{\varepsilon}{3}$。于是 $n\geqslant N$ 时
\[
  \abs{\hat f(x)-x^*_n(x)}\leqslant\abs{\hat f(x)-\hat f(y)}+\abs{\hat f(y)-x^*_n(y)}+\abs{x^*_n(y)-x^*_n(x)}\leqslant\frac\varepsilon3+\frac\varepsilon3+\frac\varepsilon3=\varepsilon .
\]

\emph{极限属于 $\Xs$}。线性函数的逐点极限是线性的，故 $\hat f$ 线性；又 $\hat f$ 是 Lipschitz 的（或者：$\abs{\hat f(x)}=\lim\abs{x^*_n(x)}\leqslant R\norm{x}$），所以连续，$\hat f\in\Xs$。
由注~\ref{rem:29}，逐点收敛即弱*收敛，故该子序列弱*收敛到 $\hat f$。
\end{proof}

\begin{theorem}[Banach--Alaoglu]
\label{res:78}
设 $\X$ 是可分赋范空间。则 $\Xs$ 的闭单位球 $B=\{x^*\in\Xs\mid\norm{x^*}_*\leqslant1\}$ 是弱*紧的。
\end{theorem}

\begin{proof}
$B$ 有界，由定理~\ref{res:76}，$B$ 上的弱*拓扑可由距离 $d$ 诱导。在度量空间中紧等价于序列紧，所以只需证明 $B$ 中任一序列 $(x^*_n)$ 都有弱*收敛到 $B$ 中某点的子序列。
由定理~\ref{res:77}，存在子序列 $(x^*_{n_k})$ 弱*收敛到某个 $x^*\in\Xs$。极限仍在 $B$ 中：对任意 $x\in\X$，
$\abs{\pair{x^*}{x}}=\lim_k\abs{\pair{x^*_{n_k}}{x}}\leqslant\norm{x}$，故 $\norm{x^*}_*\leqslant1$。
\end{proof}

\begin{example}[概率测度]
\label{exa:30}
设 $K\subseteq\R^d$ 紧，$\X=\mathcal{C}^0(K)$（由例~\ref{exa:29} 可分）。按例~\ref{exa:28}，对偶空间 $\Xs=\mathcal{M}(K)$ 是 $K$ 上的 Radon 测度空间，其单位球为
\[
  B=\{\mu\in\mathcal{M}(K)\mid\norm{\mu}_{TV}\leqslant1\},
\]
由 Banach--Alaoglu 定理，$B$ 是弱*紧的。记 $K$ 上\emph{正测度}（positive measure）的集合为
\[
  \mathcal{M}^+(K)=\{\mu\in\mathcal{M}(K)\mid\forall\phi\in\mathcal{C}^0(K),\ \phi\geqslant0\Rightarrow\pair{\mu}{\phi}\geqslant0\}.
\]
它是弱*闭半空间 $\{\mu\mid\pair{\mu}{\phi}\geqslant0\}$（$\phi\geqslant0$）之交（推论~\ref{res:74}），故弱*闭（并且是凸的）。
对 $\mu\in\mathcal{M}^+(K)$，若 $\norm{\phi}_\infty\leqslant1$，则 $\mathbf{1}_K\pm\phi\geqslant0$，从而 $\abs{\pair{\mu}{\phi}}\leqslant\pair{\mu}{\mathbf{1}_K}$，取 $\phi=\mathbf{1}_K$ 取等，
即 $\norm{\mu}_{TV}=\pair{\mu}{\mathbf{1}_K}$。所以\emph{概率测度}（probability measure）的集合
\[
  \mathcal{P}(K)=\{\mu\in\mathcal{M}^+(K)\mid\pair{\mu}{\mathbf{1}_K}=1\}=\mathcal{M}^+(K)\cap B\cap\{\mu\mid\pair{\mu}{\mathbf{1}_K}=1\}
\]
是弱*紧集 $B$ 与两个弱*闭集之交，从而是弱*紧的。
\end{example}

\subsection{自反空间}

最后把 Banach--Alaoglu 定理用于自反空间。这类空间包括 Hilbert 空间。

\begin{definition}[自反空间]
\label{def:41}
称赋范空间 $\X$ 是\emph{自反的}（reflexive），如果 $\Xss$（即 $\Xs$ 的对偶）可以等同于 $\X$，意思是典范映射
\[
  i:\X\to\Xss,\qquad x\mapsto\bigl(\phi_x:x^*\in\Xs\mapsto\pair{x^*}{x}\bigr)
\]
是双射。换言之，$\X$ 自反当且仅当 $\Xs$ 上每个连续线性形式 $\phi:\Xs\to\R$ 都能写成 $\phi(x^*)=\pair{x^*}{x}$（对某个 $x\in\X$）。
\end{definition}

典范映射总是等距的：$\norm{i(x)}=\norm{x}$（由 Hahn--Banach 定理）。本课程只需要知道关于自反空间的几个事实：
\begin{itemize}
  \item 有限维空间是自反的；
  \item Hilbert 空间是自反的；
  \item $1<p<+\infty$ 时，$\mathrm{L}^p$ 空间与 Sobolev 空间 $W^{1,p}$ 是可分且自反的；
  \item 若 $\X$ 可分且自反，则 $\Xs$ 也可分（见 [Bre10]）。
\end{itemize}

\begin{corollary}
\label{res:79}
在可分自反的赋范空间中，闭单位球是弱紧的。
\end{corollary}

\begin{proof}
记 $B_\X=\{x\in\X\mid\norm{x}\leqslant1\}$。

\emph{序列弱紧}。设 $(x_n)_{n\in\N}$ 是 $B_\X$ 中的序列。通过典范映射把它们看作 $\Xss$ 的元素：$x^{**}_n\coloneqq i(x_n)$。
由于 $\norm{i(x)}=\norm{x}$，序列 $(x^{**}_n)$ 在 $\Xss=(\Xs)^*$ 中有界。$\Xs$ 是可分赋范空间（上面最后一条事实），
因此可对 $\Xs$ 应用定理~\ref{res:77}：存在弱*收敛的子序列。重新编号后，存在 $x^{**}\in\Xss$ 使
\[
  \forall x^*\in\Xs,\quad\pair{x^{**}_n}{x^*}\xrightarrow[n\to+\infty]{}\pair{x^{**}}{x^*}.
\]
由自反性，典范映射是满射，存在 $x\in\X$ 使 $x^{**}=i(x)$。上式即
\[
  \forall x^*\in\Xs,\quad\pair{x^*}{x_n}\xrightarrow[n\to+\infty]{}\pair{x^*}{x},
\]
也就是 $(x_n)$ 弱收敛到 $x$（注~\ref{rem:27}）。$B_\X$ 是闭凸集，由命题~\ref{res:12} 是弱闭的，从而对弱收敛序列的极限封闭，故 $x\in B_\X$。

\emph{由序列紧到紧}。对 $\Xs$（可分）应用定理~\ref{res:76}：$\Xss$ 的有界集 $i(B_\X)$ 上的弱*拓扑 $\sigma(\Xss,\Xs)$ 可度量化。
而 $i$ 把 $\X$ 上弱拓扑的基本邻域 $B_{x,x^*_1,\dots,x^*_N,\varepsilon}$ 恰好映为弱*基本邻域与 $i(\X)$ 之交
\[
  B_{i(x),x^*_1,\dots,x^*_N,\varepsilon}\cap i(\X),
\]
所以 $B_\X$ 上的弱拓扑也可度量化。在可度量化空间中序列紧等价于紧，故 $B_\X$ 弱紧。
\end{proof}

\begin{remark}
上面的序列论证对任意有界闭凸集 $K$ 同样适用（把 $\norm{x_n}\leqslant1$ 换成 $\norm{x_n}\leqslant R$，最后用 $K$ 的弱闭性）：
在可分自反空间中，有界闭凸集是弱紧的（$K$ 是弱紧集 $R B_\X$ 的弱闭子集）。
\end{remark}

\begin{keypoints}
  \item 拓扑向量空间的拓扑由原点处的邻域决定（平移是同胚）；局部凸指原点有凸邻域基。
  \item 线性形式连续 $\iff$ 在原点附近有单侧界 $\iff$ 在某非空开集上有单侧界 $\iff$ 半空间 $\{\phi\leqslant\alpha\}$ 闭 $\iff$ 超平面 $\{\phi=\alpha\}$ 闭。
  \item $\sigma(\X,\Y)$ 是使 $\Y$ 中线性形式全连续的最粗拓扑，它是局部凸拓扑向量空间，且其连续线性形式恰为 $\Y$；序列收敛等价于对每个 $y\in\Y$ 有 $\pair{y}{x_n}\to\pair{y}{x}$。
  \item 弱拓扑 $\sigma(\X,\Xs)$ 比强拓扑粗：弱开（闭）$\Rightarrow$ 强开（闭）；对凸集，弱闭 $\iff$ 闭。无穷维中弱开集无界。
  \item 弱*拓扑 $\sigma(\Xs,\X)$ 下的序列收敛就是逐点收敛；它是分离的，其对偶等同于 $\X$。
  \item $\X$ 可分时，$\Xs$ 中有界集上的弱*拓扑可度量化，Banach--Alaoglu：$\Xs$ 的单位球弱*紧（例：$K$ 上的概率测度集弱*紧）。
  \item 可分自反空间中，单位球（以及有界闭凸集）弱紧。
\end{keypoints}
